\section{A non-simple symmetry family with an oscillating coefficient}%
\label{sec:asymex_czx_plus}

On a periodic chain of $N>0$ qubits, let $M$ be the bond-two tensor of
Definition~\ref{def:asymex_czx_tensors}, representing the CZX symmetry
operator of Chen, Liu, and Wen~\cite{ChenLiuWen2011CZX} (see also
Ref.~\cite[Appendix~A]{Cirac2021Matrix} and~\cite{GarreRubio2023MPOSym}),
with the first physical index denoting
the output. For a bit configuration
$t=(t_0,\ldots,t_{N-1})$, write $\bar t_n=1-t_n$ and
$E(t)=\sum_{n=0}^{N-1}t_nt_{n+1}$, with cyclic subscripts. The periodic
operator defined by these tensor entries acts as
\begin{align}
    U_N\ket{t} & =(-1)^{E(t)}\ket{\bar t}.
    \label{eq:asymex_plus_basis}
\end{align}
Thus the tensor convention used here gives $U_N=X^{\otimes N}C_N$,
where $X\ket{a}=\ket{1-a}$ and
$C_N\ket{t}=(-1)^{E(t)}\ket{t}$. For $N\geq3$, $C_N$ is the product of
nearest-neighbor controlled-$Z$ gates. The cyclic convention gives
$C_1=Z$ and $C_2=I_2$, where $Z\ket{a}=(-1)^a\ket{a}$ and
$I_N=\Id^{\otimes N}$ denotes the identity on the $N$-qubit space.
The reverse gate order $C_NX^{\otimes N}$ discussed in
Section~\ref{sec:asymex_czx} differs from this tensor convention by
the scalar $(-1)^N$; both conventions have the same square.

The basis action follows directly from the local entries
$M^{o,i}_{\alpha,\beta}
    =\delta_{o,1-i}\delta_{\beta,i}(-1)^{\alpha i}$, where all four indices
are bits and $\delta_{a,b}$ is the Kronecker delta.
In the ordinary trace contraction the virtual indices are forced to be
$g_n=t_{n-1}$, so
\begin{align}
    \tr\left(M^{s_0,t_0}\cdots M^{s_{N-1},t_{N-1}}\right)
     & =\left(\prod_{n=0}^{N-1}\delta_{s_n,1-t_n}\right)
        (-1)^{\sum_{n=0}^{N-1}t_{n-1}t_n}
    =\delta_{s,\bar t}(-1)^{E(t)}.
    \label{eq:asymex_plus_kernel}
\end{align}
Here $\delta_{s,\bar t}$ is one exactly when the two configurations
agree. This calculation also proves~(\ref{eq:asymex_plus_basis}) on
the two short rings. The flip permutes the computational basis and
its coefficients have modulus one, hence $U_N$ is unitary. Since
$E(\bar t)+E(t)=N-2\sum_n t_n+2E(t)$, its square and adjoint satisfy
\begin{align}
    U_N^2                & =(-1)^NI_N,
    \label{eq:asymex_plus_unitary}     \\
    U_N^\dagger=U_N^{-1} & =(-1)^NU_N.
    \notag
\end{align}
The same parity identity gives $C_NX^{\otimes N}=(-1)^NU_N$.

Define the operator family $O_N=I_N+U_N$. Its basis action is
\begin{align}
    O_N\ket{t} & =\ket{t}+(-1)^{E(t)}\ket{\bar t}.
    \label{eq:asymex_plus_action}
\end{align}
The identity and the CZX operator are generated by normal, inequivalent
blocks of bond dimensions one and two. Their sum therefore has a
non-simple matrix product representation of bond dimension three.
Before considering its virtual compression, its physical square follows
from~(\ref{eq:asymex_plus_unitary}):
\begin{align}
    O_N^2 & =(1+(-1)^N)I_N+2\,U_N.
    \label{eq:asymex_czx_plus_square}
\end{align}
The identity coefficient is two on even rings and zero on odd rings.

\begin{lemma}[Projectors on even rings and inverses on odd rings]
    \label{lem:asymex_plus_parity}
    For even $N>0$, $O_N/2$ is an orthogonal projector. For odd $N>0$,
    $O_N$ is invertible, with inverse $(I_N-U_N)/2$.
\end{lemma}

\begin{proof}
    For even $N$, unitarity and~(\ref{eq:asymex_plus_unitary}) give
    $U_N^\dagger=U_N$ and $U_N^2=I_N$. Consequently,
    \begin{align}
        \left(\frac{O_N}{2}\right)^\dagger & =\frac{O_N}{2},
        \notag                                                                        \\
        \left(\frac{O_N}{2}\right)^2
                                           & =\frac{I_N+2U_N+U_N^2}{4}=\frac{O_N}{2}.
        \notag
    \end{align}
    For odd $N$, $U_N^2=-I_N$. Both factors are polynomials in $U_N$, and
    \begin{align}
        (I_N+U_N)\frac{I_N-U_N}{2}
         & =\frac{I_N-U_N^2}{2}
        =I_N
        =\frac{I_N-U_N}{2}(I_N+U_N).
        \notag
    \end{align}
\end{proof}

At length zero, the unnormalized trace boundary assigns the scalar value
$U_0=[2]$ to the bond-two tensor and $O_0=[3]$ to its direct sum with the
bond-one identity, so the square law~(\ref{eq:asymex_plus_unitary}) holds only
for positive lengths.

\begin{definition}[The non-simple family and its square]
    \label{def:asymex_czx_plus}
    \lean{CZXCompression.plusTensor, CZXCompression.czxPlusIdentity,
        CZXCompression.plusBlockDim, CZXCompression.plusTargets}
    \leanok
    \uses{def:asymex_czx_tensors, def:mpdo_operator_product}
    The generating tensor is $N^{ij}=\delta^{ij}\oplus M^{ij}$ of bond
    dimension three, where $\delta^{ij}$ is the bond-one identity tensor.
    Here $N^{ij}$ denotes a local matrix, while the integer $N$ denotes
    the system size. The ordinary trace boundary gives
    \begin{align}
        \tr\left(N^{s_0,t_0}\cdots N^{s_{N-1},t_{N-1}}\right)
         & =\prod_{n=0}^{N-1}\delta_{s_n,t_n}
        +\tr\left(M^{s_0,t_0}\cdots M^{s_{N-1},t_{N-1}}\right)
        =O_N(s,t).
        \notag
    \end{align}
    The source tensor is the stacked product of bond dimension nine,
    \begin{align}
        B^{ij} & =\sum_{m=0}^{1}N^{im}\otimes N^{mj}.
        \label{eq:asymex_plus_stack}
    \end{align}
    Its four target slots, indexed by $s=0,1,2,3$, are
    $C_0=\delta$, $C_1=-\delta$, $C_2=M$, and $C_3=M$, of bond dimensions
    $1$, $1$, $2$, and $2$, respectively. Each tensor is read over the
    alphabet of physical index pairs in the order
    $(0,0)$, $(0,1)$, $(1,0)$, $(1,1)$.
\end{definition}

Stacking represents physical multiplication: expanding the intermediate
physical configuration $r$ and using multiplicativity of the Kronecker
product and $\tr(A\otimes D)=\tr(A)\tr(D)$ gives
\begin{align}
    O_N^2(s,t)
     & =\sum_r
    \tr\left(\prod_{n=0}^{N-1}N^{s_n,r_n}\right)
    \tr\left(\prod_{n=0}^{N-1}N^{r_n,t_n}\right)
    =\tr\left(\prod_{n=0}^{N-1}B^{s_n,t_n}\right).
    \label{eq:asymex_plus_contraction}
\end{align}
All matrix products are ordered by increasing site index. Thus the
periodic multiplication law is established independently of the
asymmetric fundamental theorem.

\begin{center}
    % Formula: B^{ij}=\sum_mN^{im}\otimes N^{mj}, the stacked product of
    % Definition~\ref{def:asymex_czx_plus}.
    % Source: Definition~\ref{def:asymex_czx_plus}.
    % Ink-to-index map: both rows are the bond-three generating tensor
    % $N=\delta\oplus M$; the shared vertical pairleg at each site is
    % the summed index $m$.
    % Contraction set: the vertical pairlegs (index $m$) at every site
    % and the horizontal virtual bonds within each row; the physical
    % legs $(i,j)$ remain open.
    % Boundary signature: (virtual-west=0 | virtual-east=0 |
    % physical-up=3 | physical-down=3) on both sides.
    \tenkzkernel
    \begin{tenkz}[rows={op,op}, cols=4, boundary=periodic]
        \tn[skin=mpo, ports={90:physical:$i_1$, 270:physical}]{N} &
        \tn[skin=mpo, ports={90:physical:$i_2$, 270:physical}]{N} &
        \tn[skin=dots, ports={180:virtual, 0:virtual}]{} &
        \tn[skin=mpo, ports={90:physical:$i_N$, 270:physical}]{N} \\
        \tn[skin=mpo, ports={90:physical, 270:physical:$j_1$}]{N} &
        \tn[skin=mpo, ports={90:physical, 270:physical:$j_2$}]{N} &
        \tn[skin=dots, ports={180:virtual, 0:virtual}]{} &
        \tn[skin=mpo, ports={90:physical, 270:physical:$j_N$}]{N}
    \end{tenkz}
    \quad $=$ \quad
    \begin{tenkz}[rows={op}, cols=4, boundary=periodic]
        \tn[skin=mpo, ports={90:physical:$i_1$, 270:physical:$j_1$}]{B} &
        \tn[skin=mpo, ports={90:physical:$i_2$, 270:physical:$j_2$}]{B} &
        \tn[skin=dots, ports={180:virtual, 0:virtual}]{} &
        \tn[skin=mpo, ports={90:physical:$i_N$, 270:physical:$j_N$}]{B}
    \end{tenkz}
\end{center}

\begin{theorem}[Compression of the squared non-simple family]
    \label{thm:asymex_czx_plus_compression}
    % Principal declaration: CZXCompression.czxPlusIdentity_isReduction.
    \lean{CZXCompression.czxPlusIdentity_compression,
        CZXCompression.czxPlusIdentity_isReduction,
        CZXCompression.czxPlusIdentity_left_mul_right_of_ne,
        CZXCompression.czxPlusIdentity_dim_eq,
        CZXCompression.czxPlusIdentity_evalWord_remainder_eq_zero}
    \leanok
    \uses{def:asymex_czx_plus, thm:asym_multiblock}
    The source of Definition~\ref{def:asymex_czx_plus} carries the block
    triangular form of Theorem~\ref{thm:asym_multiblock} for the four
    slots with $z=3$ zero slots and $9=1+1+2+2+3$. Writing $d_s$ for
    the dimension of $C_s$, the compression maps
    $W_s:\C^9\to\C^{d_s}$ and $V_s:\C^{d_s}\to\C^9$ satisfy
    \begin{align}
        W_sV_t & =
        \begin{cases}
            \Id_{d_s}, & s=t,    \\
            0,         & s\ne t.
        \end{cases}
        \notag
    \end{align}
    For every word $w$ over the alphabet of physical index pairs,
    $W_sB^wV_s=C_s^w$. The remainder
    $R^{ij}=B^{ij}-\sum_{s=0}^3V_sC_s^{ij}W_s$ satisfies $R^w=0$
    whenever $|w|\geq7$.
\end{theorem}

\begin{proof}
    \leanok
    \uses{thm:asymex_explicit_gauge}
    The splitting $\C^3=\C\oplus\C^2$ induces a four-summand
    decomposition of $\C^3\otimes\C^3$. After grouping these summands,
    each stacked letter has the form
    \begin{align}
        B^{ij} & \cong
        \delta^{ij}\oplus M^{ij}\oplus M^{ij}
        \oplus\left(\sum_{m=0}^1M^{im}\otimes M^{mj}\right).
        \notag
    \end{align}
    The fourth summand is the squared tensor of
    Section~\ref{sec:asymex_czx}. Its flag contributes the target
    $-\delta$ and three zero slots. The other three summands already
    carry $\delta$, $M$, and $M$.

    Assemble the change of bond coordinates from these four summands.
    Checking block triangularity and the matched and unmatched diagonal
    blocks is a finite verification over integer matrices.
    Theorem~\ref{thm:asymex_explicit_gauge} then gives the compression
    pairs. The dimension count and the nilpotency bound are clauses
    (vii) and (vi) of Theorem~\ref{thm:asym_multiblock}, with
    $|S|+z=4+3=7$.
\end{proof}

\begin{theorem}[The oscillating structure coefficient]
    \label{thm:asymex_czx_plus_trace}
    \lean{CZXCompression.czxPlusIdentity_trace_evalWord}
    \leanok
    \uses{def:asymex_czx_plus}
    For every nonempty word $w$ over the alphabet of index pairs,
    \begin{align}
        \tr(B^w) & =(1+(-1)^{|w|})\tr(\delta^w)+2\tr(M^w).
        \label{eq:asymex_czx_plus_trace}
    \end{align}
\end{theorem}

\begin{proof}
    \leanok
    \uses{thm:asymex_czx_plus_compression, prop:asym_compression_trace}
    Proposition~\ref{prop:asym_compression_trace} applied to
    Theorem~\ref{thm:asymex_czx_plus_compression} gives, for nonempty $w$,
    \begin{align}
        \tr(B^w)
         & =\sum_{s=0}^3\tr(C_s^w)
        =\tr(\delta^w)+\tr((-\delta)^w)+2\tr(M^w).
        \notag
    \end{align}
    Each letter of the second slot contributes a factor $-1$, so
    $\tr((-\delta)^w)=(-1)^{|w|}\tr(\delta^w)$.
    This proves~(\ref{eq:asymex_czx_plus_trace}).
\end{proof}

Taking $|w|=N>0$, (\ref{eq:asymex_plus_contraction}) identifies
(\ref{eq:asymex_czx_plus_trace}) with the physical square
law~(\ref{eq:asymex_czx_plus_square}). The compression resolves the
oscillating integer coefficient into the signed weights $+1$ and $-1$:
the coefficient $1+(-1)^N$ is their length-$N$ power sum, as in
Theorem~\ref{thm:asym_unblocked_coefficients}. The two copies of $M$
each have unit weight and contribute the constant coefficient two.
This is the general power-sum mechanism, not an instance of the positive
diagonal matrices in the CPSV characterization of MPDO renormalization
fixed points. In particular, the even-ring projector property does not
establish an MPDO renormalization-fixed-point condition.
